\clearpage
\section{Sidequest: quantum mechanics over the field with one element}
\label{sec:sidequest-f1}

\textbf{Opened 6 September 2026.} This is a separate, literature-grounded
exploration of arithmetic quantum mechanics over $\mathbb F_1$, also written
$\mathbb F_{\mathrm{un}}$. Its north star is a precisely formulated quantum
theory, with particular attention to the analogue of a qubit and to composition.
The different established analogies are retained as separate candidates.
The accompanying research map is \path{docs/sidequests/f1-qm.md}; primary
source bodies, versions and retrieval hashes are indexed in
\path{refs/LEDGER.md}. The finite-ring campaign is not superseded.

There is substantial relevant literature, but the constructions preserve
different structures. A field may specify the arithmetic base, the state
coefficients, a counting parameter, or a deformation parameter. Confusing
these roles would make the comparison misleading. In the present campaign,
the arithmetic base is a finite field or ring and the amplitudes remain
complex. The direct ``quantum $\mathbb F_{\mathrm{un}}$'' papers instead
change the coefficient object of the states themselves.

\subsection{What could count as a qubit?}

Two distinguishable states, a two-dimensional Hilbert space, a simple object
of categorical dimension two, and two simple objects of a category are
different notions. The following table makes the distinction explicit.

\begin{longtable}{@{}p{3.1cm}p{5.1cm}p{6.5cm}@{}}
\toprule
Approach & Qubit analogue & Composition \\
\midrule
\endhead
Pointed $\mathbb F_1$ modules & A basepoint and two basis points. Complex
linearization of normal partial maps gives $(\mathbb C^2,M_2(\mathbb C))$.
& Smash product has four nonzero basis points and realizes as the Hilbert
tensor product. Native basis states have no superpositions.\\[3pt]
Cyclotomic modules & Two free $\mu_N$-orbits. For configuration $C_2$ and
phase level two, clock and shift give the Heisenberg qubit.
& Balanced smash of modules; central product of Heisenberg groups; ordinary
tensor after complex realization.\\[3pt]
Absolute Cartesian frames & Rank two has $N+2$ projective rays. At bare
$\mathbb F_1$ these are $10,01,11$.
& The rank-four frame has fifteen rays at $N=1$: nine product rays and six
nonproduct rays. The operation on pairs is the Kronecker rule.\\[3pt]
Thin symplectic geometry & Two coordinate polarizations, suggestive of
position and momentum. These are not two Hilbert-state vectors.
& Orthogonal sums of phase geometry; Hilbert tensor requires a realization.\\[3pt]
Blueprint Heisenberg geometry & A chosen $\mathbb F_2$ fibre gives the
reference order-eight group and its two-dimensional central-character module.
& Ring realization and phase data precede the physical tensor product.\\[3pt]
Bands and crowds & A relational Heisenberg object, with 27 signed or eight
Krasner points. Neither number is a count of qubit states.
& Geometric products exist; physical composition needs a specified phase
and representation construction.\\[3pt]
Hall/groupoid oscillator & No preferred qubit in one oscillator. Two colours
give a two-dimensional one-particle sector after realization.
& Full Fock spaces tensor on disjoint colour sets. Two identical bosons in
two modes have dimension three; two distinguishable qubits have dimension four.\\[3pt]
Representation categories & The two-dimensional simple object of
$\operatorname{Rep}(D_8)$ or $\operatorname{Rep}(Q_8)$.
& Internal fusion splits its square into four charges. Independent copies
instead use an external tensor and a matching-central-character quotient.\\[3pt]
Quantum torus & $VU=-UV$, $U^2=V^2=1$ gives $M_2(\mathbb C)$.
& Two selected central fibres give $M_4(\mathbb C)$.\\[3pt]
Arithmetic descent and Bost--Connes & A specified finite fibre or a chosen
two-level encoding; no canonical qubit is supplied by the global arithmetic data.
& Frobenius and transfer correspondences; tensor products of statistical
systems after realization.\\[3pt]
Tropical and motivic routes & Support or incidence information, without a
specified positive two-state quantum system.
& The relevant semiring, geometric or categorical tensor must be named.\\
\bottomrule
\end{longtable}

\subsection{The developed geometric and physical analogies}

\paragraph{Counting and thin geometry.}
The projective count $1+q+\cdots+q^{r-1}$ specializes to $r$ at $q=1$;
Gaussian coefficients specialize to subset counts. Tits's thin geometry
therefore retains coordinate points, subsets and permutation symmetries.
For symplectic rank $r$ the reflection group is
$W(C_r)=(C_2)^r\rtimes S_r$, acting on symplectic pairs. A coordinate
Lagrangian selects one direction from each pair, giving $2^r$ choices.
This is the incidence content of the familiar specialization of
$\prod_{i=1}^r(q^i+1)$.
The geometric analogies are developed in
\href{https://arxiv.org/abs/math/0404185}{Deitmar's \emph{Schemes over
$\mathbb F_1$}, section 5}, and
\href{https://arxiv.org/abs/1201.1324}{Lorscheid's Tits--Weyl models}.
Here a \emph{Weyl group} is a reflection group, not the group of Weyl
displacement operators. Moreover, $q^3\mapsto1$ for the order of
$\mathrm{UT}_3(\mathbb F_q)$ is a polynomial specialization, not a limit
of groups or a proof that its $\mathbb F_1$ geometry is empty.

\paragraph{Direct quantum $\mathbb F_{\mathrm{un}}$.}
\href{https://arxiv.org/abs/1312.4191}{Chang, Lewis, Minic and Takeuchi}
construct a $q=1$ coordinate-state specialization of finite-field amplitude
models: superposition disappears. This is different from this campaign's
complex Hilbert space $\ell^2(\mathbb F_q)$.
\href{https://arxiv.org/abs/1808.09694}{Thas's \emph{Absolute Quantum Theory}}
uses larger Cartesian frames and develops monomial dynamics and
quantum-information analogies. Its Hermitian expression is partial: a sum
is allowed only when at most one term is nonzero. Orthogonality is
disjointness of supports. The involution determines the unitary subgroup;
inversion of phases makes all phase-permutation matrices unitary in this
sense. Its cloning discussion distinguishes simple projective rays, which
can be copied, from all frame states. These statements do not supply an
ordinary positive-definite Hilbert inner product on the frame.

\paragraph{An existing geometric Heisenberg model.}
Lorscheid's Proposition 5.5 constructs Tits--Weyl models for standard
parabolic unipotent radicals of $\mathrm{GL}_n$. For the upper-triangular
Borel of $\mathrm{GL}_3$, this is the reference Heisenberg group
\[
h(a,b,t)=\begin{pmatrix}1&a&t\\0&1&b\\0&0&1\end{pmatrix},\qquad
h(a,b,t)h(c,d,u)=h(a+c,b+d,t+u+ad).
\]
Its underlying blue scheme is affine three-space. Its Weyl extension has
one point, while ring base extension recovers the full unipotent group.
Thus a positive $\mathbb F_1$ model already exists; selecting a central
character and constructing its representation are further steps. The
whole polarizing-cocycle torsor, especially its characteristic-two types,
is not supplied by that reference model. Likewise, a Tits lift of a
reflection is not automatically a metaplectic or Weil lift.

\paragraph{Cyclotomic conventions.}
For the raw phase monoid $\{0\}\cup\mu_N$, integral realization gives
$\mathbb Z[\mu_N]\simeq\mathbb Z[t]/(t^N-1)$. A cyclotomic blueprint may
instead impose additive relations giving $\mathbb Z[\zeta_N]$.
A faithful complex phase character is another named realization datum.
These conventions are compared in
\href{https://arxiv.org/abs/1301.0083}{Lorscheid's \emph{A blueprinted view
on $\mathbb F_1$-geometry}}, section 1.1.4. The group $\mu_{q-1}$ reflects
multiplicative field data; finite-field Weyl phases come from additive
characters and lie in $\mu_p$. These two phase counts must not be identified.

\subsection{Pointed modules, ordinary qubits and absolute frames}

\begin{definition}[Normal pointed maps and their two products]
For $r\geq0$ let $V_r=\{*,1,\ldots,r\}$. A normal pointed map preserves
$*$ and has at most one preimage of each non-basepoint; it is a partial
injection of the nonzero sets. Wedge identifies basepoints, while smash
collapses all pairs containing a basepoint in the Cartesian product.
Put $L(S)=\mathbb C[S\setminus\{*\}]$, with its basis orthonormal, and
let $B(S)$ be the complex span of realized normal endomorphisms in
$\operatorname{End}_{\mathbb C}L(S)$.
\end{definition}
\begin{scope}These objects have no native addition of vectors. The image
algebra $B(S)$ is explicitly a complex realization.\end{scope}
\provenance{D1010}{definitions.md}{f1\_check.py (M13)}{supporting comparison}

\begin{proposition}[The pointed-set qubit after operator realization]
{\normalfont\statusSketch}\quad For every positive-rank finite pointed set $S$,
$B(S)=\operatorname{End}_{\mathbb C}L(S)$. For finite pointed sets $S,T$,
\[
L(S\wedge T)\simeq L(S)\otimes L(T).
\]
For positive ranks this identifies the image algebras with their tensor
product. In particular, $V_2$ realizes to $(\mathbb C^2,M_2(\mathbb C))$,
and $V_2\wedge V_2=V_4$ realizes to two qubits.
\end{proposition}
\begin{scope}The realization retains a distinguished basis. The concrete
image algebra is not the unreduced partial-injection monoid algebra.
Algebras are transported along bijections; arbitrary partial maps do not
give unital maps between full endomorphism algebras.\end{scope}
\provenance{F1-NORMAL}{f1-comparisons.md, section 3}{f1\_check.py (M13)}{supporting SKETCH}
\begin{proof}
The partial map taking $j$ to $i$ and everything else to the basepoint
realizes as $E_{ij}$. These matrix units span the full algebra, and their
partial inverses realize adjoints. Smash products of these maps give
$E_{ij}\otimes E_{k\ell}$. The basis identification proves the tensor claim.
\end{proof}

At rank two the partial-injection monoid has seven elements, including
zero. Its contracted monoid algebra has dimension six, whereas its image
has dimension four: linear realization imposes relations such as
$I=E_{00}+E_{11}$. Keeping only the two bijections gives the commutative
span of $I$ and the swap. Thus the class of allowed morphisms matters.
The pointed-module, normal-map and base-change conventions come from
\href{https://arxiv.org/abs/1204.5395}{Szczesny}, pages 5--8.

\begin{definition}[Cyclotomic projective Cartesian frames]
Fix an abstract cyclic group $\mu_N$ of order $N\geq1$ and put
$F_N=\{0\}\sqcup\mu_N$, with absorbing zero. For $r\geq1$, define
\[
P_N(r)=\bigl(F_N^r\setminus\{(0,\ldots,0)\}\bigr)/\mu_N
\]
by simultaneous scalar multiplication. The product-state map
$P_N(r)\times P_N(s)\to P_N(rs)$ sends $([x],[y])$ to
$[(x_i y_j)_{i,j}]$. Rays outside its image are called nonproduct rays
in this factorization sense.
\end{definition}
\begin{scope}A Cartesian frame is different from a free pointed module.
No probability rule or total Hermitian form is stipulated.\end{scope}
\provenance{D1009}{definitions.md}{f1\_check.py (M11)}{supporting comparison}

\begin{proposition}[Frame qubits and their nonproduct rays]
{\normalfont\statusSketch}\quad For every $N\geq1$,
\[
|P_N(r)|=\frac{(N+1)^r-1}{N}.
\]
The rank-two product-state map is injective, so the rank-four frame has
$(N+2)^2$ product rays and $N(N+1)(N+2)$ nonproduct rays.
\end{proposition}
\begin{scope}Nonproduct means failure of tensor factorization. It does not
by itself establish a Born rule, interference or Bell correlations.\end{scope}
\provenance{F1-FRAME}{f1-comparisons.md, section 2}{f1\_check.py (M11)}{supporting SKETCH}
\begin{proof}
Scalar multiplication acts freely on any nonzero tuple. For a nonzero
product array, ratios within a nonzero row and column recover both
factors up to scalars. The ray counts then follow by subtraction.
\end{proof}

For $N=1$ the one-system rays are $10,01,11$; the composite has fifteen
rays, nine product and six nonproduct. The diagonal array
$\left(\begin{smallmatrix}1&0\\0&1\end{smallmatrix}\right)$ is a nonproduct
because product supports are rectangular. This larger frame is the
reason a nonadditive theory can retain a factorization notion of
entanglement while the coordinate-only model does not. In the strict
pointed-map category, only simple vectors are morphisms from the unit;
the additional frame states require an explicitly enlarged state notion.

\subsection{A finite cyclotomic Weyl--Heisenberg formulation}

\begin{definition}[Cyclotomic modules and named complex realization]
Fix $N\geq1$, an abstract cyclic group $\mu_N$ of order $N$, and a named
faithful character $\iota:\mu_N\to\mathbb C^\times$. Let
$\mathrm{Free}_*^{\mu_N}$ be finite pointed sets with a $\mu_N$ action
fixing the basepoint and free on its complement; maps are pointed
equivariant maps. Define
\[
\mathbb C_\iota(S)=\mathbb C[S]/(e_0,\ e_{us}-\iota(u)e_s)
\]
and use the induced linear maps on morphisms.
\end{definition}
\begin{scope}The phase embedding is part of the realization datum. The
objects are free phase sets, not Cartesian frames.\end{scope}
\provenance{D1001}{definitions.md}{f1\_check.py}{f1-r1.md}

\begin{definition}[Hyperbolic phase data and the pointed Schr\"odinger action]
For a finite abelian group $A$ with exponent dividing $N$, put
\[
A^\vee=\operatorname{Hom}(A,\mu_N),\quad V_A=A\times A^\vee,
\quad c_A((a,\chi),(b,\eta))=\eta(a)^{-1},
\]
\[
\kappa_A((a,\chi),(b,\eta))=\chi(b)\eta(a)^{-1}.
\]
Define $H_N(A)=\mu_N\times V_A$ by
$(u,v)(u',w)=(uu'c_A(v,w),v+w)$ and adjoin an absorbing zero to get
$H_N(A)^0$. Put $S_A=\{0\}\sqcup(\mu_N\times A)$ and stipulate
\[
(u,a,\chi)[z,x]=[uz\chi(x+a),x+a],\qquad 0_Hs=0_S.
\]
On its realization let $X(a)e_x=e_{x+a}$,
$Z(\chi)e_x=\iota(\chi(x))e_x$, and $W(a,\chi)=Z(\chi)X(a)$.
\end{definition}
\begin{scope}The abelian configuration group is extra data. No field
structure, division by two or identification $A\simeq A^\vee$ is assumed.\end{scope}
\provenance{D1002,D1003}{definitions.md}{f1\_check.py (M1--M3)}{f1-r1.md}

\begin{theorem}[The finite cyclotomic Schr\"odinger system]
{\normalfont\statusProved}\quad Fix $N\geq1$, cyclic $\mu_N$ of order $N$,
a faithful $\iota:\mu_N\to\mathbb C^\times$, and a finite abelian $A$
of exponent dividing $N$. For the associated phase system,
$|A^\vee|=|A|$, characters extend from subgroups and separate points, and
$\kappa_A$ is perfect. The displayed multiplication makes $H_N(A)$ a
group with centre exactly $\mu_N$, acting faithfully on $S_A$.
After the named complex realization,
\[
W(a,\chi)W(b,\eta)=\iota(\eta(a)^{-1})W(a+b,\chi\eta),\qquad
\operatorname{Tr}(W(v)^*W(w))=|A|\delta_{v,w}.
\]
These operators span $\operatorname{End}_{\mathbb C}(\mathbb C[A])$.
The irreducible unitary representation of $H_N(A)$ with central character
$\iota$ is unique up to equivalence, has dimension $|A|$, and has
intertwiners unique up to $U(1)$.
\end{theorem}
\begin{scope}This is finite Heisenberg theory for abelian configurations,
with an explicit pointed module. A universal representation theorem for
all proposed $\mathbb F_1$ geometries is not asserted.\end{scope}
\provenance{F1-DUAL,F1-WEYL,F1-REAL}{f1-cyclotomic.md, sections 0--2}{f1\_check.py (M1--M3)}{f1-adjudication.md}
\begin{proof}
Character extension follows by adjoining a cyclic generator and choosing
a root in $\mu_N$; the exponent hypothesis ensures the required root exists.
Restriction and induction on the group order give $|A^\vee|=|A|$.
A nontrivial character sums to zero by translating its sum. Character
evaluation verifies the cocycle identity and action. Testing the two
coordinate subgroups proves perfection. The trace vanishes unless both
translation and character are trivial; the resulting $|A|^2$ orthogonal
operators form a basis of the full matrix algebra. Matrix units then prove
uniqueness of its simple module and scalar uniqueness of intertwiners.
\end{proof}

\begin{definition}[Fixed-level isomorphisms and balanced composition]
Let $\mathrm{FinAb}_N^{\simeq}$ have the above configurations and their
group isomorphisms. For $f:A\to B$, put
$f_*\chi=\chi\circ f^{-1}$,
$S(f)[u,x]=[u,f(x)]$ and
$H(f)(u,a,\chi)=(u,f(a),f_*\chi)$.
Define
\[
S\wedge_{\mu_N}T=(S\wedge T)/([us,t]\sim[s,ut]),\qquad
H\boxdot K=(H\times K)/\{(u,u^{-1}):u\in\mu_N\}.
\]
In the last expression the two groups have specified central copies of
$\mu_N$.
\end{definition}
\begin{scope}Functoriality is asserted on isomorphisms at a common phase
level. Arbitrary configuration homomorphisms are not silently included.\end{scope}
\provenance{D1004}{definitions.md}{f1\_check.py (M6)}{f1-r1.md}

\begin{theorem}[Composition of independent phase systems]
{\normalfont\statusProved}\quad Fix $N\geq1$ and a faithful phase realization
$\iota$. The module, group and Hilbert assignments on
$\mathrm{FinAb}_N^{\simeq}$ are strong symmetric monoidal, with comparisons
\begin{gather*}
S_A\wedge_{\mu_N}S_B\simeq S_{A\times B},\qquad
H_N(A)\boxdot H_N(B)\simeq H_N(A\times B),\\
\mathbb C[A]\otimes\mathbb C[B]\simeq\mathbb C[A\times B].
\end{gather*}
At $N=1$ the exponent condition forces $A=0$ and a one-dimensional
phase-system realization, while free pointed modules of arbitrary rank
still exist.
\end{theorem}
\begin{scope}The group tensor is a central quotient, not a Cartesian
product retaining two independent centres. The level-one statement is
restricted to perfect phase systems, not all $\mathbb F_1$ quantum routes.\end{scope}
\provenance{F1-FUNCT,F1-ONE}{f1-cyclotomic.md, sections 3,5}{f1\_check.py (M1,M6)}{f1-adjudication.md}
\begin{proof}
The phase-set comparison sends $[[u,a],[v,b]]$ to $[uv,(a,b)]$.
The corresponding group map multiplies central phases and takes the
external character. Its kernel is the antidiagonal centre. Every
associativity diagram multiplies the same phases and associates the
same tuple, proving coherence. The Hilbert comparison is the basis map.
\end{proof}

\begin{definition}[Central pushout from a finite local ring]
Let $R$ be finite commutative unital local with $1\ne0$, and let
$\psi:(R,+)\to\mu_N$ be named, with exponent of $(R,+)$ dividing $N$.
Assume $\iota\psi$ is generating, meaning its kernel contains no nonzero
ideal. Let $H_{\beta_0}(R)$ have triples $(t,a,b)$ and product
$(t+u+ad,a+c,b+d)$. Put
\[
\chi_b(x)=\psi(-bx),\qquad
P_\psi=\bigl(\mu_N\times H_{\beta_0}(R)\bigr)
/\{(\psi(t)^{-1},(t,0,0)):t\in R\}.
\]
\end{definition}
\begin{scope}A generating character need not be faithful on the additive
group. The raw Heisenberg group and its phase image are distinguished.\end{scope}
\provenance{D1005,D12,D13,D16}{definitions.md}{f1\_check.py (M4,M5)}{f1-r1.md}

\begin{theorem}[Exact recovery of the arithmetic reference convention]
{\normalfont\statusProved}\quad Let $R$ be a finite commutative unital local
ring with $1\ne0$, exponent of $(R,+)$ dividing $N\geq1$, and a named
$\psi:(R,+)\to\mu_N$ such that the faithful complex realization
$\iota\psi$ is generating. Then
$b\mapsto\chi_b$ identifies $(R,+)$ with its phase dual and
$P_\psi\simeq H_N((R,+))$ by
$[u;(t,a,b)]\mapsto(u\psi(t),a,\chi_b)$.
The raw group map has central kernel $\ker\psi$. Its realized operators are
\[
W(a,\chi_b)e_y=(\iota\psi)(-b(y+a))e_{y+a}=Z(-b)X(a)e_y,
\]
with product cocycle $(\iota\psi)(ab')$, uniformly including characteristic two.
\end{theorem}
\begin{scope}Only the reference cocycle is recovered here. This does not
identify all raw polarizing-cocycle groups with one phase group.\end{scope}
\provenance{F1-RING}{f1-cyclotomic.md, section 4}{f1\_check.py (M4,M5)}{f1-adjudication.md}

For example, over $\mathbb F_4$ a nontrivial trace character has a
two-element kernel. Its raw Heisenberg group has order 64 and its phase
image order 32. Over $\mathbb F_2$, the order-eight reference group acts
on the usual qubit. This supplies a checked bridge to the main campaign.

\subsection{Fourier transform and a precisely scoped additive extension}

\begin{definition}[Quadratic phases, strict normalizer and Fourier kernel]
A quadratic phase is $\theta:A\to\mu_N$, $\theta(0)=1$, for which
$B_\theta(a,x)=\theta(a+x)\theta(a)^{-1}\theta(x)^{-1}$ is a bicharacter.
For $t\in A$, $f\in\operatorname{Aut}(A)$ put
$T_{t,f,\theta}[u,x]=[u\theta(x),f(x)+t]$. The strict normalizer is the
normalizer of the $\mu_N$-level Weyl group in
$\operatorname{Aut}_{\mathrm{Free}_*^{\mu_N}}(S_A)$. Define
\[
F_Ae_x=\sum_{\rho\in A^\vee}\iota(\rho(x))e_\rho,\quad
\operatorname{ev}_a(\rho)=\rho(a),\quad
J_A(a,\chi)=(\chi^{-1},\operatorname{ev}_a),\quad r_A(a,\chi)=\chi(a).
\]
\end{definition}
\begin{scope}A self-Fourier endomorphism additionally needs a pairing
identifying $A$ with its dual. No such pairing is suppressed.\end{scope}
\provenance{D1006}{definitions.md}{f1\_check.py (M7)}{f1-r1.md}

\begin{theorem}[What remains monomial and what Fourier requires]
{\normalfont\statusProved}\quad Fix $N\geq1$, a faithful $\iota$, and finite
abelian $A$ of exponent dividing $N$. Its strict normalizer consists exactly of the
$T_{t,f,\theta}$ and central phases. Its phase-space maps are
\[
(a,\chi)\longmapsto
\bigl(f(a),(\chi B_\theta(a,-))\circ f^{-1}\bigr).
\]
For $|A|>1$, $F_A$ has full column support and cannot be the realization
of a strict pointed map. Nevertheless,
\[
F_A^*F_A=|A|I,\qquad
F_AW(v)=\iota(r_A(v))W(J_Av)F_A.
\]
\end{theorem}
\begin{scope}Only bicharacters admitting a $\mu_N$-valued quadratic
refinement occur in the strict subgroup. Enlarging phases alone does not
make Fourier monomial. No canonical Weil splitting is asserted.\end{scope}
\provenance{F1-MON}{f1-cyclotomic.md, section 6}{f1\_check.py (M7, Fourier only)}{f1-adjudication.md}

\begin{definition}[Framed cyclotomic kernels and lifted phase isomorphisms]
Put $K_N=\mathbb Q[\mu_N]/\ker\iota$. The category $\operatorname{Mat}(K_N)$
has finite sets $I$ as objects and $J\times I$ matrices over $K_N$ as
morphisms; tensor is Cartesian product of sets and Kronecker product of
matrices. The associated modules are explicitly framed by $[1,i]$ in
$S_I=\{0\}\sqcup(\mu_N\times I)$. Let $\operatorname{PMat}(K_N)$ be
invertible matrices modulo $K_N^\times$, realized by entrywise evaluation.

The groupoid $\operatorname{LiftHyp}_N$ has configurations $A$ as objects.
An arrow is $(g,r)$ with $g:V_A\to V_B$ a group isomorphism preserving
$\kappa$ and $r:V_A\to\mu_N$ satisfying
\[
r(v)r(w)c_B(gv,gw)=c_A(v,w)r(v+w).
\]
It denotes the centre-fixing map $(u,v)\mapsto(ur(v),g(v))$; composition
composes these maps. Tensor is $A\times B$ on objects and
$(g\times h,(v,w)\mapsto r(v)s(w))$ on arrows, with the canonical
reordering of phase factors. The unit is $A=0$ and the symmetry is the swap.
\end{definition}
\begin{scope}This target is an additive cyclotomic envelope. It is not
claimed to be a category of $\mathbb F_1$ schemes; arbitrary free torsors
are not given coordinate matrices without a frame.\end{scope}
\provenance{D1007}{definitions.md}{no general checker}{f1-adjudication.md, O1--O3 repair}

\begin{proposition}[Projective realization of lifted phase symmetries]
{\normalfont\statusSketch}\quad For $N\geq1$ and a named faithful $\iota$,
there is a projective strong symmetric monoidal functor
\[
Q_N:\operatorname{LiftHyp}_N\longrightarrow\operatorname{PMat}(K_N)
\]
defined by the Egorov intertwining equations. Its complex realization is
projectively unitary; affine-quadratic maps and Fourier have the values above.
\end{proposition}
\begin{scope}The lift is input data. Geometric descent of the additive
envelope and a lift of every symplectic transformation remain separate questions.\end{scope}
\provenance{F1-CORR}{f1-cyclotomic.md, section 7}{not claimed as computationally proved}{supporting SKETCH after scope repair}
\begin{proof}
The intertwining equations have coefficients in $K_N$ and a one-dimensional
solution space after embedding in $\mathbb C$. Gaussian elimination has
the same rank over either field, so there is a nonzero solution over
$K_N$. Irreducibility makes it invertible; uniqueness up to scalar gives
projective composition and tensor coherence. Unitary normalization is
performed only after complex realization.
\end{proof}

\subsection{Relational Heisenberg geometry over bands}

\begin{definition}[Bands and the upper-unitriangular crowd]
A band is a commutative pointed multiplicative monoid $B$ with a null
ideal $N_B$ of formal sums, in which every $a$ has a unique $-a$ with
$a+(-a)\in N_B$; morphisms preserve null sums. The regular partial field
$\mathbb F_1^\pm=\{0,1,-1\}$ uses sums vanishing in $\mathbb Z$.
The Krasner hyperfield $\mathbb K=\{0,1\}$ declares a sum null exactly
when it has zero or at least two nonzero terms.
Define $H_{\mathrm{crowd}}(B)=B^3$, represented by $h(a,b,t)$ as above,
with identity $(0,0,0)$. Its colaw consists of triples for which
\[
\sum_i a_i,\qquad \sum_i b_i,\qquad
t_1+t_2+t_3+a_1b_2+a_1b_3+a_2b_3
\]
and both cyclic variants of the last sum lie in $N_B$.
The inverse set of $h$ consists of $k$ with $(h,k,1)$ in the colaw;
the product set $h*k$ consists of $c$ for which some $d$ has both
$(h,k,d)$ and $(c,d,1)$ in the colaw.
\end{definition}
\begin{scope}Nullity is part of the band, not an invented binary addition.
The signed band and the raw pointed phase monoid have different structure.\end{scope}
\provenance{D1008}{definitions.md}{f1\_check.py (M10)}{supporting comparison}

\begin{proposition}[A band-level Heisenberg crowd and its ring fibres]
{\normalfont\statusSketch}\quad The functor $B\mapsto H_{\mathrm{crowd}}(B)$
on bands, with the displayed colaw, is an affine algebraic crowd.
For every commutative ring $R$, its fibre is exactly $\mathrm{UT}_3(R)$.
The signed fibre has 27 points, 319 colaw triples and 23 points admitting
an inverse inside that fibre. The Krasner fibre has eight points and
152 colaw triples, with multivalued inverses.
\end{proposition}
\begin{scope}These finite fibres are not ordinary groups. Even a crowd
product with the identity may be empty when its required inverse mediator
is absent. No Hilbert representation is inferred from the point counts.\end{scope}
\provenance{F1-CROWD}{f1-comparisons.md, section 1}{f1\_check.py (M10)}{supporting SKETCH}

This specializes the algebraic $\mathrm{SL}_3$ crowd of
\href{https://arxiv.org/abs/2305.13809}{Lorscheid--Thas (2023)}, section 5.4.
Restricting to an identity-containing subset preserves the crowd axioms;
the displayed coordinate relations are affine. Over rings they say
$h_1h_2h_3=I$. The signed example retains
$h(1,0,0)*h(0,1,0)=\{h(1,1,1)\}$, and hence the central cross-term,
although $h(1,0,0)^2$ is absent. This provides a constructive second
Heisenberg formulation, with a representation theory still to specify.

\subsection{Tensor categories and the qubit fusion object}

\begin{definition}[The categorical phase system]
Let $\mathcal C_N(A)=\operatorname{Rep}_{\mathbb C}(H_N(A))$, with all
linear intertwiners and internal tensor using the diagonal group action.
Let $\mathcal C_N(A)_\lambda$ be the full subcategory on which central
$u$ acts as $\lambda(u)I$. The unitary version chooses invariant Hermitian
inner products and has the same intertwining spaces. Write $\omega_A$
for the forgetful tensor functor. For groups $H,K$ with specified central
$\mu_N$, define
$\operatorname{Rep}(H)\boxtimes_{\mathrm{match}}\operatorname{Rep}(K)$
as the full subcategory of $\operatorname{Rep}(H\times K)$ on which all
$(u,u^{-1})$ act trivially.
\end{definition}
\begin{scope}The matching condition is explicit. It is not an
unspecified relative tensor product, and it is different from internal fusion.\end{scope}
\provenance{D1011}{definitions.md}{f1\_check.py (examples)}{supporting comparison}

\begin{proposition}[Central sectors explain why the category matters]
{\normalfont\statusSketch}\quad For $N\geq1$ and finite abelian $A$ of
exponent dividing $N$, $\mathcal C_N(A)$ is a symmetric fusion
category with a faithful tensor functor $\omega_A$, and
\[
\mathcal C_N(A)=\bigoplus_{\lambda\in\widehat{\mu_N}}
\mathcal C_N(A)_\lambda,\qquad
\mathcal C_N(A)_\lambda\otimes\mathcal C_N(A)_\nu
\subseteq\mathcal C_N(A)_{\lambda\nu}.
\]
Moreover, for finite groups $H,K$ with specified central embeddings of $\mu_N$,
\[
\operatorname{Rep}(H\boxdot K)\simeq
\operatorname{Rep}(H)\boxtimes_{\mathrm{match}}\operatorname{Rep}(K).
\]
On simple external tensors, the matching condition is equality of the
two central characters. In particular a fixed nontrivial sector is not
a tensor category by itself.
\end{proposition}
\begin{scope}The unit of matching composition on this family is
$\operatorname{Rep}(\mu_N)$, with one simple per grade. It is not the
unit for the unrestricted Deligne product.\end{scope}
\provenance{F1-CAT}{f1-comparisons.md, section 5}{f1\_check.py (M6,M12; examples)}{supporting SKETCH}
\begin{proof}
Average an inner product over the finite group to obtain semisimplicity.
Central character eigenspaces give the grading, and diagonal tensor
multiplies eigencharacters. A representation factors through the central
quotient precisely when its antidiagonal centre acts trivially; inflation
and descent are inverse tensor functors. On an external simple the
condition is $\lambda(u)\nu(u)^{-1}=1$.
\end{proof}

\begin{definition}[Dihedral and quaternion qubit data]
For $e\in\{0,1\}$ let $G_e=\mathbb F_2^3$ with
\[
(t,a,b)(u,c,d)=(t+u+ad+e(ac+bd),a+c,b+d).
\]
Put
\[
\sigma_e(t,a,b)=(-1)^t i^{e(a+b)}Z^bX^a,\quad
X=\begin{pmatrix}0&1\\1&0\end{pmatrix},\quad
Z=\begin{pmatrix}1&0\\0&-1\end{pmatrix},
\]
using representatives $a,b\in\{0,1\}$ in the exponent of $i$. Define
$\nu_2(\sigma_e)=|G_e|^{-1}\sum_g\operatorname{Tr}(\sigma_e(g^2))$.
\end{definition}
\begin{scope}The group law distinguishes the two cocycle choices. The
quaternion matrices require fourth-root phases in this monomial realization.\end{scope}
\provenance{D1012}{definitions.md}{f1\_check.py (M12)}{supporting comparison}

\begin{proposition}[Two qubit objects with the same fusion rules]
{\normalfont\statusSketch}\quad $G_0=D_8$ of order eight and $G_1=Q_8$. Each has four
one-dimensional characters and the single two-dimensional irreducible
$\sigma_e$. Writing the four characters as $1,\alpha,\beta,\alpha\beta$,
\[
\sigma_e\otimes\sigma_e\simeq
1\oplus\alpha\oplus\beta\oplus\alpha\beta,\qquad
\nu_2(\sigma_0)=1,\quad\nu_2(\sigma_1)=-1.
\]
Their sourced tensor-category identifications are
\[
\operatorname{Rep}(D_8)\simeq\mathrm{TY}(\mathbb F_2^2,\chi_{\rm alt},1/2),
\qquad
\operatorname{Rep}(Q_8)\simeq\mathrm{TY}(\mathbb F_2^2,\chi_{\rm alt},-1/2).
\]
\end{proposition}
\begin{scope}The displayed tensor square uses one diagonal symmetry
group. It does not replace the external tensor describing independent qubits.\end{scope}
\provenance{F1-FUSION}{f1-comparisons.md, section 6; Shimizu, p. 20}{f1\_check.py (M12)}{supporting SKETCH}

The two-dimensional character is $2,-2$ on the two central elements and
zero elsewhere. Squaring it gives the sum of the four quotient characters.
Six elements square to $1$ in $D_8$ and only two do in $Q_8$, giving the
indicator signs. The categorical identifications and their distinction
are in \href{https://arxiv.org/abs/1005.4500}{Shimizu's
\emph{Frobenius--Schur indicators in Tambara--Yamagami categories}},
Definition 3.1, Table 1 and page 20.

For two independent reference qubits the relevant group instead is
$H_2(C_2)\boxdot H_2(C_2)=H_2(C_2\times C_2)$, of order 32.
Its nontrivial-central-character irreducible has dimension four.
Both this representation and internal fusion use the underlying Hilbert
space $\mathbb C^2\otimes\mathbb C^2$, but their symmetry actions differ.
This is a concrete reason to retain the tensor category, its sectors and
its fibre functor as part of the quantum object.

The fibre functor is essential to this particular formulation:
\[
\operatorname{End}_{\operatorname{Rep}(H)}(\sigma)=\mathbb C,\qquad
\operatorname{Hom}_{\operatorname{Rep}(H)}(1,\sigma)=0
\]
for the nontrivial central character. The equivariant category alone does not give every
state vector or every qubit observable. Those are recovered as vectors
of $\omega(\sigma)=\mathbb C^2$ and elements of
$\operatorname{End}_{\mathbb C}(\omega(\sigma))=M_2(\mathbb C)$.
An alternative categorical state notion would have to be specified.

\subsection{The independent Hall and categorical-oscillator route}

\begin{definition}[Pointed-set Hall algebra and Fock realization]
On the normal pointed-set category above, the rational Hall vector space
has basis $u_r$ indexed by rank. Its product counts pointed subsets $T$
with the prescribed isomorphism classes of $T$ and $S/T$.
Put $\mathcal F_{\rm alg}=\mathbb C[x]$ with
$\langle x^m,x^n\rangle=\delta_{mn}n!$, define
$a^\dagger f=xf$, $af=df/dx$, and let $\mathcal F_{\rm bos}$ be the
Hilbert completion.
\end{definition}
\begin{scope}Hall coefficients and the Fock space are additive
realizations of the pointed-set data. Creation is unbounded.\end{scope}
\provenance{D1010}{definitions.md}{f1\_check.py (M8,M9)}{supporting comparison}

\begin{proposition}[A bosonic Heisenberg algebra from finite sets]
{\normalfont\statusSketch}\quad The rational Hall algebra of finite normal
pointed sets is $\mathbb Q[x]$ under
\[
u_mu_n=\binom{m+n}{m}u_{m+n},\qquad u_n\longmapsto x^n/n!.
\]
On $\mathcal F_{\rm alg}$, creation and annihilation are adjoints and
$[a,a^\dagger]=I$. In the normalized basis $x^n/\sqrt{n!}$ their
coefficients are $\sqrt{n+1}$ and $\sqrt n$.
\end{proposition}
\begin{scope}The polynomial Hall algebra alone is not the Weyl algebra;
annihilation is additional operator structure. No finite-dimensional
characteristic-zero CCR representation is asserted.\end{scope}
\provenance{F1-HALL}{f1-comparisons.md, section 4}{f1\_check.py (M8,M9)}{supporting SKETCH}

The subset count proves the Hall formula. The operator calculation is
$d(xf)/dx-x\,df/dx=f$. The source construction is
\href{https://arxiv.org/abs/1204.5395}{Szczesny's Hall algebra}, section 3.
\href{https://arxiv.org/abs/math/0004133}{Baez--Dolan} give the factorial
groupoid weights and the positive categorical relation
\[
A A^\dagger\simeq A^\dagger A\oplus\mathrm{id}.
\]
\href{https://arxiv.org/abs/1207.2054}{Morton--Vicary} realize the
categorified Heisenberg algebra using spans of finite-set groupoids:
add/remove histories split according to whether the removed point was
just added. Their two-linearization retains the tower of
$\operatorname{Rep}(S_n)$, with the stated finiteness/completion caveats.
The connection to $\mathbb F_1$ is the equivalence between finite sets and
the core of finite pointed modules. It is not a claim that these authors
construct one universal $\mathbb F_1$ scheme theory.

This supplies a categorical north star different from the finite Weyl
system. Degree two of the representation-category tower has two simples
in $\operatorname{Rep}(S_2)$; that is not automatically a qubit. With
two colours, the one-particle complex space is a qubit and the two-boson
sector is $\operatorname{Sym}^2(\mathbb C^2)$, of dimension three.
A one-mode cutoff to $|0\rangle,|1\rangle$ instead has projected
commutator $\operatorname{diag}(1,-1)$, and is not closed under creation.
The Hall operation uses wedge and particle-number addition; the finite
configuration tensor uses smash and rank multiplication.

\subsection{Toric, analytic and arithmetic-statistical formulations}

\begin{definition}[A quantum torus and its selected central fibre]
For a named $\xi\in\mathbb C^\times$ put
\[
\mathcal T_\xi=\mathbb C\langle U^{\pm1},V^{\pm1}\rangle/(VU-\xi UV).
\]
If $\xi$ is a primitive $N$th root, put
$\mathcal T_\xi(1,1)=\mathcal T_\xi/(U^N-1,V^N-1)$.
\end{definition}
\begin{scope}The parameter $\xi$, its order $N$ and a finite-field
cardinality are different data. Central values are explicitly selected.\end{scope}
\provenance{D1013}{definitions.md}{cyclic Weyl checks}{supporting comparison}

\begin{proposition}[The qubit as a quantum-torus fibre]
{\normalfont\statusSketch}\quad For $\xi$ primitive of order $N$,
$\mathcal T_\xi(1,1)\simeq M_N(\mathbb C)$ via
$Ue_j=e_{j+1}$ and $Ve_j=\xi^je_j$ on $\mathbb Z/N$.
For $N=2$ this is a qubit algebra; two independent fibres tensor to
$M_4(\mathbb C)$.
\end{proposition}
\begin{scope}The universal quantum torus has not been identified with
a finite matrix algebra. At $\xi=1$ it is commutative.\end{scope}
\provenance{F1-TORUS}{f1-comparisons.md, section 7}{f1\_check.py (cyclic cases)}{supporting SKETCH}

The $N^2$ clock/shift words are trace-orthogonal and span, which proves
the selected-fibre statement. This is the twisted-lattice picture of
\href{https://arxiv.org/abs/math/0511263}{Neeb's rational quantum tori}.
Monoidal tori retain character lattices; a chosen skew pairing is
additional symplectic data. \href{https://arxiv.org/abs/0809.1564}{Manin's
cyclotomic analytic geometry} supplies Habiro completions and evaluation
or Taylor maps at roots of unity. It suggests an arithmetic interpolation
problem for quantum fibres, but does not itself give a common Hilbert
space or compatible Weil intertwiners. A deformation limit $\xi\to1$
is different from all the preceding $q\to1$ operations.

\paragraph{Frobenius descent.}
\href{https://arxiv.org/abs/0906.3146}{Borger's lambda geometry} treats
lambda structure as descent data below $\operatorname{Spec}\mathbb Z$.
For flat schemes it can be expressed as commuting Frobenius lifts.
Monoid algebras have the lift $m\mapsto m^p$. Its naive coordinatewise
application to the integral Heisenberg group fails already on the
additive subgroup: $(1+1)^p\ne1^p+1^p$ in $\mathbb Z$.
This rejects that particular group-compatible lift, not all arithmetic
descent. A quantum extension should specify whether arithmetic maps act
by homomorphisms, noninvertible operators or correspondences. No qubit
or Born rule is selected by lambda structure alone.

\paragraph{A published quantum-statistical realization.}
\href{https://arxiv.org/abs/0806.2401}{Connes--Consani--Marcolli's
\emph{Fun with $\mathbb F_1$}}, Proposition 6.1 and Theorem 6.2,
constructs an $\mathbb F_1$ model of the Bost--Connes endomotive, with
the rational system recovered by scalar extension. It uses a cyclotomic
tower and Frobenius/transfer correspondences. Its complex realization
has the time evolution
\[
\sigma_t(\mu_n)=n^{it}\mu_n,\qquad \sigma_t(e(r))=e(r).
\]
In the standard representation on $\ell^2(\mathbb N_{>0})$,
$He_n=(\log n)e_n$ and
$\operatorname{Tr}(e^{-\beta H})=\zeta(\beta)$ for $\beta>1$;
see \href{https://repo-archives.ihes.fr/FONDS_IHES/I_Prepublications/CONNES/1994-1998/M_95_38/M_95_38_web.pdf}{Bost--Connes (1995)}.
This is genuine quantum statistical dynamics, not a finite
Stone--von Neumann representation. Choosing two energy levels provides
an encoding but not a sector invariant under its full observable algebra.

\href{https://arxiv.org/abs/1901.00020}{Lieber--Manin--Marcolli}, published
in 2022, develop lifts to Grothendieck rings, assembler categories,
spectra and Nori motives. This is an independent literature-based reason
to explore categorical targets, without asserting that these categories
are the fusion categories of the finite Heisenberg construction.

\paragraph{Tropical and motivic boundaries.}
The band/Krasner framework retains combinatorial flags beyond a single
coordinate apartment. Boolean semirings, Krasner hyperfields, signed
partial fields and pointed monoids have different additive information.
\href{https://arxiv.org/abs/1405.4527}{Connes--Consani's Arithmetic Site}
uses a tropical semiring on a multiplicative-semigroup topos to realize
arithmetic correspondences. It does not prescribe a local Weyl algebra
or positive quantum states. Relative and homotopical geometries broaden
the candidate bases and targets, but no general Heisenberg representation
theorem for those bases is established in this sidequest.
Finally, \href{https://arxiv.org/abs/1209.4837}{Bejleri--Marcolli's
\emph{Quantum field theory over $\mathbb F_1$}} concerns torifications
and varieties attached to Feynman integrals; its introduction explicitly
distinguishes that problem from physical Lagrangians and Feynman rules
over $\mathbb F_1$.

\paragraph{An additional tensor analogy: fermionic occupancy.}
\href{https://arxiv.org/abs/1307.4522}{Lin--Wang--Wu--Yang (2013)}
categorify the one-mode fermion algebra with
$c^2=(c^\dagger)^2=0$ and $cc^\dagger+c^\dagger c=1$.
Occupancy subsets suggest a combinatorial exterior-power route; this
is our proposed $\mathbb F_1$ bridge, not their claim.
\href{https://doi.org/10.4230/LIPIcs.FSCD.2018.17}{Hadzihasanovic--de
Felice--Ng (2018)}, section 2, give a precise monoidal category of local
fermionic modes and maps of definite parity. Its basic object has
underlying space $\mathbb C^{1|1}$; tensor adds parity and the physical
fermionic swap negates $|11\rangle$. The two occupancy labels therefore
do not supply unrestricted qubit operations. Parity-preserving one-mode
observables are diagonal, while a fixed-parity two-mode sector can encode
a qubit. Signed phases are a candidate carrier for exchange signs over
$\mathbb F_1$; the cited complex categorical construction does not prove
that descent. This is a further composition law to compare with the
bosonic and distinguishable-system routes.

\subsection{Precisely formulated outputs and open comparisons}

The finite arithmetic candidate is the configuration-to-phase-system
functor, with a named phase realization and a checked tensor law. The
categorical candidate retains
\[
\bigl(\mathcal C_N(A),\ \mathcal C_N(A)_\iota,\ \omega_A\bigr)
\]
together with internal fusion and matching external composition. The
pointed normal-map candidate already realizes an ordinary qubit over
the bare combinatorial base. The Heisenberg crowd is a geometric object
with correct ring fibres, while the Hall/groupoid route gives an
independent categorical oscillator. Bost--Connes supplies the arithmetic
quantum-statistical test case. These are several precise formulations,
not evidence that all the analogies must coincide.

The next comparison problems are concrete: relate a chosen Tits reflection
lift to its projective Fourier operator; extend the geometric Heisenberg
model to nonreference cocycles; identify a genuine geometric home for the
additive phase kernels; compare finite cyclotomic fibres with arithmetic
transfer maps; and distinguish tensor products of distinguishable systems
from bosonic particle-number sectors. A theory on arbitrary $\mathbb F_1$
schemes still needs its input category, morphisms and realization functor
specified. It is a research target, not a theorem or a well-formed
universal conjecture supplied by the present literature.

Seven finite-kernel claims are \statusProved{}; the extensions remain
\statusSketch{}. Thirteen exact example gates and their mutations check
the finite computations. The cited literature and our deductions are
distinguished throughout.
